\documentclass[11pt]{article}
\usepackage[T1]{fontenc}
\usepackage{lmodern}
\usepackage[margin=1in]{geometry}
\usepackage{amsmath,amssymb,amsthm,booktabs}
\usepackage[hidelinks]{hyperref}
\newtheorem{theorem}{Theorem}
\newtheorem{lemma}[theorem]{Lemma}
\newcommand{\Z}{\mathbb Z}
\title{Extremal degree multiplicities for modified total\newline
irregularity of chemical graphs}
\author{}
\date{October 13, 2026}
\begin{document}
\maketitle

\begin{abstract}
For a connected simple graph on $n$ vertices with maximum degree at most
four, consider the sum of $|d(u)-d(v)|^{1/n}$ over all unordered pairs of
distinct vertices. We determine every degree-multiplicity vector attaining
the maximum, for every $n\geq7$. This proves the classification proposed in
Conjecture~10 of Knor, \v Skrekovski, Filipovski and Dimitrov. The proof first
optimizes over integer multiplicities satisfying the handshaking parity
condition. A positive quadratic decomposition reduces the problem to two
nearest-integer comparisons. We then show that every maximizing vector has
a connected simple realization. The argument applies to all $n\geq7$ and
does not depend on finite enumeration.
\end{abstract}

\section{Statement and relation to previous work}

A \emph{chemical graph} here means a finite, undirected, connected simple
graph of maximum degree at most four, as in~\cite{knor}. For such a graph
$G$ of order $n\geq7$, define
\begin{equation}\label{eq:index}
 I_n(G)=\sum_{\{u,v\}\in\binom{V(G)}{2}}|d(u)-d(v)|^{1/n}.
\end{equation}
The sum is over all unordered pairs of distinct vertices, not only edges.
Write $a_i=|\{v\in V(G):d(v)=i\}|$ for $1\leq i\leq4$.

In arXiv:2411.01530v1, Knor et al.~\cite[Theorem~9 and Conjecture~10]{knor}
proved a classification
for a smaller, order-dependent range of positive exponents, and conjectured
that the same classification holds for exponent $1/n$.
Dimitrov et al.~\cite[Theorem~2.2]{dimitrov} subsequently reported the
$n\equiv0\pmod4$ case, within a more general parameter statement.
We give a self-contained proof for all four residue classes; no result
about the more general exponent $c/n$ is used below.
The constant-exponent Conjecture~11 of~\cite{knor} is a different question.

\begin{theorem}\label{thm:main}
For every integer $n\geq7$, a chemical graph of order $n$ maximizes
\eqref{eq:index} if and only if its degree-multiplicity vector occurs in
Table~\ref{tab:winners}. Each listed vector has a connected simple
realization.
\begin{table}[ht]
\centering
\begin{tabular}{cl}
\toprule
Order & $(a_1,a_2,a_3,a_4)$\\
\midrule
$4k$ & $(k,k,k,k)$\\
$4k+1$ & $(k,k,k,k+1)$\\
$4k+2$ & $(k,k+1,k,k+1)$ or $(k+1,k,k+1,k)$\\
$4k+3$ & $(k+1,k,k+1,k+1)$\\
\bottomrule
\end{tabular}
\caption{All maximizing degree multiplicities, with $n\geq7$.}
\label{tab:winners}
\end{table}
\end{theorem}

The last row is the vector $(K,K-1,K,K)$ when $n=4K-1$, as written
in~\cite{knor}. The classification concerns degree multiplicities, not
isomorphism types. Any two graphs with the same multiplicities have the
same value of~\eqref{eq:index}.

\section{Optimization over integer multiplicities}

Every chemical graph gives a vector in the larger set
\[
 \mathcal A_n=\left\{a\in\Z_{\geq0}^4:
 \sum_{i=1}^4a_i=n,\quad a_1+a_3\equiv0\pmod2\right\}.
\]
The parity condition follows from the handshaking lemma. We first maximize
over $\mathcal A_n$ without requiring graphicality. Put
\[
 p=\frac1n,\qquad T=2^p,\qquad S=3^p,\qquad D=1+2T-S.
\]
Counting pairs by their degrees gives
\begin{equation}\label{eq:Q}
 Q(a)=a_1a_2+a_2a_3+a_3a_4+T(a_1a_3+a_2a_4)+Sa_1a_4.
\end{equation}
For every graph with multiplicities $a$, its index equals $Q(a)$.

Set $s=a_1+a_4$, $u=a_1-a_4$, and $v=a_2-a_3$. The inverse change of
variables is
\begin{equation}\label{eq:inverse}
 a_1=\frac{s+u}{2},\quad a_4=\frac{s-u}{2},\quad
 a_2=\frac{n-s+v}{2},\quad a_3=\frac{n-s-v}{2}.
\end{equation}
Expansion of~\eqref{eq:Q} yields
\begin{align}
 4Q(a)&=n^2+2nTs-Ds^2-E(u,v),\label{eq:decomposition}\\
 E(u,v)&=Su^2+v^2+2(T-1)uv\nonumber\\
 &=(S-T+1)u^2+(2-T)v^2+(T-1)(u+v)^2.\label{eq:positive}
\end{align}
All three coefficients in~\eqref{eq:positive} are strictly positive:
$1<T<2$ and $S>T$. The integer and parity conditions become
\begin{equation}\label{eq:congruences}
 u\equiv s\pmod2,\qquad v\equiv n-s\pmod2,\qquad
 u-v\equiv-n\pmod4.
\end{equation}
Nonnegativity additionally requires $0\leq s\leq n$, $|u|\leq s$,
and $|v|\leq n-s$.

\subsection{The parity penalty}

The congruences give the bounds and equality conditions in
Table~\ref{tab:penalty}. These are uniform lower bounds; at an arbitrary
boundary value of $s$, their equality conditions need not be feasible.
\begin{table}[ht]
\centering
\begin{tabular}{ccccc}
\toprule
 & \multicolumn{2}{c}{$s$ even} & \multicolumn{2}{c}{$s$ odd}\\
$n\bmod4$ & Bound on $E$ & Equality $(u,v)$ & Bound on $E$ & Equality $(u,v)$\\
\midrule
0 & $0$ & $(0,0)$ & $S+2T-1$ & $(1,1),\ (-1,-1)$\\
1 & $1$ & $(0,1)$ & $S$ & $(-1,0)$\\
2 & $4$ & $(0,2),\ (0,-2)$ & $S+3-2T$ & $(1,-1),\ (-1,1)$\\
3 & $1$ & $(0,-1)$ & $S$ & $(1,0)$\\
\bottomrule
\end{tabular}
\caption{Lower bounds for $E(u,v)$ under~\eqref{eq:congruences}.}
\label{tab:penalty}
\end{table}

Here is a verification of every case. If $n\equiv0\pmod4$ and $s$ is
even, positivity in~\eqref{eq:positive} gives $E\geq0$, with equality
only at $(0,0)$. If $n\equiv0\pmod4$ and $s$ is odd, then $u,v$ are
odd and $u+v\equiv2\pmod4$. Thus the three squares in~\eqref{eq:positive}
are at least $1,1,4$, respectively. Equality requires $u=v=\pm1$.

If $n$ is odd and $s$ is even, then $u$ is even, $v$ is odd, and $u+v$
is odd. Consequently $E\geq(2-T)+(T-1)=1$. Equality requires
$u=0$, $|v|=1$, and the last congruence in~\eqref{eq:congruences}
fixes the sign of $v$. If $n$ and $s$ are odd, the corresponding bound is
$E\geq(S-T+1)+(T-1)=S$, with $|u|=1$, $v=0$ and the sign of $u$
again fixed by that congruence.

If $n\equiv2\pmod4$ and $s$ is odd, then $u,v$ are odd, giving
$E\geq(S-T+1)+(2-T)=S+3-2T$. Equality holds exactly when
$u=-v=\pm1$. Finally, if $n\equiv2\pmod4$ and $s$ is even,
then $u,v$ are even, $u-v\equiv2\pmod4$, and $u+v\equiv2\pmod4$.
Hence $u^2+v^2\geq4$ and $(u+v)^2\geq4$. Since
$S-T+1>2-T$, we obtain
\[
 E\geq4(2-T)+4(T-1)=4,
\]
with equality exactly at $u=0$, $v=\pm2$.

\subsection{The nearest integers and comparison of the two parities}

Define
\[
 H(s)=\frac{n^2+2nTs-Ds^2}{4}.
\]
Strict concavity of $x^p$ on the positive real axis gives
$2T>1+S$, so $D>2$. The vertex of $H$ is
\begin{equation}\label{eq:vertex}
 s_* = \frac{nT}{D}=\frac n2+\delta,
 \qquad \delta=\frac{n(S-1)}{2D}.
\end{equation}
The binomial theorem, with $n\geq2$, gives
$(1+2/n)^n>1+n(2/n)=3$. Thus
\begin{equation}\label{eq:bounds}
 n(S-1)<2,\qquad 0<\delta<\frac12.
\end{equation}

For fixed $n$, let $L(s)$ denote the lower bound from
Table~\ref{tab:penalty} for the parity of $s$. Then
$Q(a)\leq U(s):=H(s)-L(s)/4$.
On each parity class $L(s)$ is constant; consequently its upper bound
$U(s)$ is uniquely maximized at the integer of that parity nearest to
$s_*$. Equation~\eqref{eq:bounds} makes these two integers explicit in
Table~\ref{tab:comparison}.
\begin{table}[ht]
\centering
\begin{tabular}{ccccc}
\toprule
$n\bmod4$ & Even candidate & Odd candidate & Winning $s$ & Winning $U$ minus other $U$\\
\midrule
0 & $n/2$ & $n/2+1$ & $n/2$ & $(4T-n(S-1))/4$\\
1 & $(n-1)/2$ & $(n+1)/2$ & $(n+1)/2$ & $(n-1)(S-1)/4$\\
2 & $n/2+1$ & $n/2$ & $n/2$ & $(2D-n(S-1))/4$\\
3 & $(n+1)/2$ & $(n-1)/2$ & $(n+1)/2$ & $(n+1)(S-1)/4$\\
\bottomrule
\end{tabular}
\caption{The two parity candidates and their strict comparison.}
\label{tab:comparison}
\end{table}

Both candidates have $s\geq3$ and $n-s\geq3$ when $n\geq7$.
Every equality pair in the corresponding row of
Table~\ref{tab:penalty} therefore satisfies the nonnegativity constraints
in~\eqref{eq:inverse}. Thus both upper bounds being compared are attained
in $\mathcal A_n$. Substitution into $U$ gives the differences in the last
column. Each is strictly positive by $S>1$, $T>1$, $D>2$, and
\eqref{eq:bounds}.

The winning $(s,u,v)$ are therefore
\[
\begin{array}{c|c}
n & (s,u,v)\\ \hline
4k & (2k,0,0)\\
4k+1 & (2k+1,-1,0)\\
4k+2 & (2k+1,1,-1)\ \text{or}\ (2k+1,-1,1)\\
4k+3 & (2k+2,0,-1).
\end{array}
\]
Using~\eqref{eq:inverse} gives exactly Table~\ref{tab:winners}.
The strict comparison between parities, uniqueness of the nearest integer
within each parity, and the equality conditions for $E$ exclude every
other maximizing vector in $\mathcal A_n$.

\section{Connected simple realizations}

It remains to show that each optimal integer vector is realized by a
chemical graph. We first give a small-degree graphicality lemma.

\begin{lemma}\label{lem:graphical}
If $n\geq8$ and $d_1,\ldots,d_n\in\{1,2,3,4\}$ have even sum, then
they are the degrees of a simple graph.
\end{lemma}
\begin{proof}
Arrange the sequence in nonincreasing order. By the Erd\H os--Gallai
criterion~\cite{erdosgallai}, it suffices to check, for $1\leq r\leq n$,
\[
 \sum_{i=1}^r d_i\leq r(r-1)+\sum_{i=r+1}^n\min(r,d_i).
\]
For $r=1$, the left side is at most $4$ and the right side is $n-1\geq7$.
For $r=2$, the respective bounds are $8$ and $2+(n-2)=n\geq8$.
For $r=3$, they are $12$ and $6+(n-3)=n+3$, which suffice for $n\geq9$.
If $n=8$ and the inequality at $r=3$ fails, integrality forces the left
side to be $12$ and the right side to be $11$. The degree sequence must
then be $(4,4,4,1,1,1,1,1)$, whose sum is odd, a contradiction.
For $r=4$, the left side is at most $16$ and the right side at least
$12+(n-4)=n+8\geq16$. For $r\geq5$, the term $r(r-1)$ alone is at
least $4r$, which bounds the left side.
\end{proof}

For completeness, we also prove the connected-realization criterion
recorded in~\cite[(1.6)]{erdosgallai}.
\begin{lemma}\label{lem:connected}
A graphical sequence of $n$ positive integers with sum at least
$2(n-1)$ has a connected simple realization.
\end{lemma}
\begin{proof}
Start with any simple realization. If it has $c\geq2$ components and
every component is acyclic, then it has $n-c\leq n-2$ edges, contrary to
the assumed degree sum. Thus a component contains a cycle. Choose an
edge $uv$ of such a cycle and an edge $xy$ in another component, the
latter existing because all degrees are positive. Replace $uv,xy$ by
$ux,vy$. The new edges did not previously exist and have distinct
endpoints, so the graph remains simple and every degree is unchanged.
Removing $uv$ leaves its component connected. If $xy$ was a bridge,
its two resulting pieces both join the first component through the new
edges; otherwise its component remains connected and joins the first.
In either case the number of components decreases by one. Repeating
this operation gives a connected realization.
\end{proof}

\begin{proof}[Completion of the proof of Theorem~\ref{thm:main}]
The vectors in Table~\ref{tab:winners} have even degree sums, respectively
\[
 10k,\qquad 10k+4,\qquad 10k+6\ \text{or}\ 10k+4,
 \qquad 10k+8.
\]
Each sum is at least $2(n-1)$ in the stated range $n\geq7$.
For $n\geq8$, Lemmas~\ref{lem:graphical} and~\ref{lem:connected} therefore
give connected simple realizations with all degrees at most four.
For $n=7$, the required sequence is $(4,4,3,3,2,1,1)$; on vertices
$0,\ldots,6$ the edge set
\[
 \{02,03,04,05,12,13,14,16,23\}
\]
is an explicit connected simple realization, where $ij$ denotes $\{i,j\}$.
Thus the maximum over $\mathcal A_n$ is attained in the original graph
class. Since every chemical graph gives a vector in $\mathcal A_n$, the
classification already obtained proves the theorem.
\end{proof}

\section{Reproducible finite checks}

The supplied \texttt{verify.py} script uses only the Python standard
library. It enumerates all $317{,}826$ parity-feasible multiplicity vectors for
$7\leq n\leq60$. It encloses $2^{1/n}$ and $3^{1/n}$ between rational
bounds with denominator $2^{48}$, certifying the bounds by integer powers.
For every nonmaximizing vector it verifies a strictly positive lower
bound on the difference from the stated maximum. The two maximizing
vectors when $n\equiv2\pmod4$ have identical integer coefficient triples
in the basis $(1,T,S)$, so their equality is exact.
The script also constructs and checks all $242$ listed maximizing
degree sequences for $7\leq n\leq200$, using degree-preserving edge
switches to obtain connected realizations. Edge lists and numerical
certificates are included with the script. These are finite consistency
checks, not premises of the proof for arbitrary $n$.

\paragraph{AI use and verification status.}
This research was completed with AI assistance. The mathematical argument
and reproducible checks are supplied for scrutiny. This version does not
claim verification in a proof assistant or external peer review.

\begin{thebibliography}{9}
\bibitem{knor}
M.~Knor, R.~\v Skrekovski, S.~Filipovski, and D.~Dimitrov,
\emph{Extremizing antiregular graphs by modifying total $\sigma$-irregularity},
Applied Mathematics and Computation \textbf{490} (2025), 129199.
\href{https://doi.org/10.1016/j.amc.2024.129199}{doi:10.1016/j.amc.2024.129199}.
Preprint \href{https://arxiv.org/abs/2411.01530v1}{arXiv:2411.01530v1}
(2024), especially Theorem~9 and Conjecture~10.

\bibitem{dimitrov}
D.~Dimitrov, S.~Filipovski, M.~Knor, and R.~\v Skrekovski,
\emph{Modified Total Sigma-Irregularity in Chemical Graphs},
Discrete Mathematics Letters \textbf{18} (2026), 1--7.
\href{https://doi.org/10.47443/dml.2026.041}{doi:10.47443/dml.2026.041}.

\bibitem{erdosgallai}
P.~Erd\H os and T.~Gallai,
\emph{Gr\'afok el\H o\'irt fok\'u pontokkal},
Matematikai Lapok \textbf{11} (1960), 264--274 (in Hungarian).
\href{https://users.renyi.hu/~p_erdos/1961-05.pdf}{Full-text archive}.
\end{thebibliography}
\end{document}
